\honest{Focus Your Uncertainty}{Eliezer Yudkowsky}{2007}

Will bond yields go up, go down or stay the same? If you are a TV pundit, it does not matter: you can explain any outcome afterwards. But suppose you are a novice, with 100 minutes to prepare your excuses in advance. Now you would like to know which outcome will happen.

You cannot use ``probabilities,'' of course. Probabilities are ``little numbers'' that appear next to word problems at school, and there are none here. Besides, you feel uncertain. I say all this in your voice, and I do not mean it.

You cannot spend all 100 minutes on each of the three outcomes. A review committee would make you split the time equally, since nothing written down justifies any other split, but there is no committee. There is a Federal Reserve announcement tomorrow, so ``remain the same'' seems unlikely.

Plausibility does not help, because all three outcomes are plausible, and your ability to explain them is unlimited. Your time is not. And yet you do expect the outcomes differently, and when one seems more likely, the others seem less likely. That is ``the fascinating part.''

The relation between how much you expect each outcome and how much time you give it ``can't actually be quantified''; then a parenthesis adds that your payoff is ``logarithmic in time spent preparing the excuse.'' \nb{That aside is what quantifies it. With a logarithmic payoff, the best plan gives each excuse minutes in proportion to its probability. With a payoff that grows in proportion to time, all 100 minutes go to the likeliest outcome.} So perhaps anticipation is a conserved resource, like money, to be allocated.

Your statistics courses never taught you this. If you used numbers at all, you might end up with ``pairs of numbers'' that you call ``DS for Dexter-Sinister,'' though you do have only 100 minutes. \nb{``DS'' is presumably Dempster--Shafer theory, which gives each hypothesis a lower and an upper number. It also divides a fixed total, among sets of outcomes, so the fixed 100 minutes do not by themselves rule it out. The post offers only the pun.}

If only there were an art of focusing your uncertainty. What could we call it? \nb{The answer implied is probability theory applied to degrees of belief, worked out by F. P. Ramsey in 1926 and L. J. Savage in 1954, among others. The post names neither.}
